\documentclass{ximera}
\title{Linear Functions}


\newcommand{\pskip}{\vskip 0.1 in}

\begin{document}
\begin{abstract}
Review of linear functions in the context of motion and rotation.
\end{abstract}
\maketitle


\section{Motion Along a Line}
\emph{Position} and \emph{velocity} are vectors and we'll learn more about vectors later in the class. But for now we'll confine our attention to motion along a straight line where position and velocity are scalars. Put a one-dimensional coordinate system on this line and you have a number line, say the $s$-axis. We can now regard the position of a beetle on the number line as \emph{scalar} (ie. a number). It's just the $s$-coordinate of the beetle and it might be positive or negative. The $s$-coordinate measures position of the beetle relative to the origin, $s=0$.

\emph{Velocity} is the rate of change of position with respect to time. For our beetle crawling along a line, velocity is also a scalar. It is positive when the beetle is moving in the positive direction of the number line (ie. when its $s$-coordinate increases). The beetle's velocity is negative when its $s$-coordinate decreases.

We'll start with a special case, where the beetle's velocity is \emph{constant}. Here's an example.

\begin{example} \label{Ex:lfle0311}
Between noon and 1:00pm a beetle crawls along the $s$-axis with a constant velocity, passing position $s=5$cm at 12:23pm and position $s=-12$cm at 12:30pm.

\begin{onlineOnly}
    \begin{center}
\desmos{ddf6exycnx}{900}{600}
\end{center}
\end{onlineOnly}

\href{https://www.desmos.com/calculator/ddf6exycnx}{142: Beetle on a Line}


\begin{enumerate}
\item Find the beetle's position at 12:51pm.

\item Find a function
\[
      s= f(t) \, , \, 0\leq t \leq 60,
\] 
that expresses the beetle's position (in cm) in terms of the number of minutes past noon.

\end{enumerate}

\begin{explanation}
\begin{enumerate}
\item We start by finding the beetle's velocity. Since the velocity is constant, the velocity is the beetle's change in position divided by the change in time. In this example, the beetle's velocity is
\[
  v = \frac{\Delta s}{\Delta t} = \frac{(-12 - 5) \text{ cm}}{(30-23)\text{ min}} = \frac{-17}{7} \, \frac{\text{cm}}{\text{min}}. 
\]

To find the beetle's position at 12:51pm, we first find its change in position during the time from 12:23pm to 12:51pm. To do this we multiply its velocity by the number of minutes $51 - 23 = 28$ between these times. So the change in position is
\begin{align*}
      f(51) - f(23) &= \left(  \frac{-17}{7} \, \frac{\text{cm}}{\text{min}} \right) \left( 51 - 23\right) \text{min} \\
                        &= \left(  \frac{-17}{7} \, \frac{\text{cm}}{\text{min}} \right) \left( 28 \text{ min} \right)  \\
                        &= -68 \text{ cm}.
\end{align*}

Now to get the beetle's position at 12:51pm we add this change in position to its position at 12:23pm. This gives the beetle's position at 12:51pm to be
\begin{align*}
    f(51) &= f(23) + (f(51) - f(23))\\
            &= 5 \text{ cm} + \left(  \frac{-17}{7} \, \frac{\text{cm}}{\text{min}} \right) \left( 28 \text{ min} \right)  \\
            &= 5 \text{ cm} + -68 \text{ cm} \\
            &= -63 \text{ cm}
\end{align*}

\item To find the beetle's position at time $t$ minutes past noon, we use the same idea as in part (a). The only difference is that we replace the time $51$ minutes past noon with the time $t$ minutes past noon. So the beetle's position (in cm) at time $t$ minutes past noon is
\begin{align*}
     s &= f(t) \\
        &= f(23) + (f(t) -  f(23)) \\
        &=  5 \text{ cm} + \left(  \frac{-17}{7} \, \frac{\text{cm}}{\text{min}} \right) \left( t - 23 \right) \text{min} \\
        &=  5 \text{ cm} + \left(  \frac{-17}{7} \right) \left(t - 23 \right) \text{ cm} .
\end{align*}

Usually we supress the units, include the domain, and write the position function as
\[
    s = f(t) =  5 - \frac{17}{7}\left(t - 23  \right) \, , \, 0\leq t \leq 60.
\]

Explain the logic behind this expression. Include units in your explanation.
\begin{freeResponse}
\end{freeResponse}

\end{enumerate}
\end{explanation}

\end{example}

\section{Uniform Circular Motion}

Uniform circular motion is the basis of much of trigonometry. Uniform here means constant speed. So a beetle moving in uniform circular motion crawls around a circle at a constant speed.

For a circle in an $xy$-coordinate system, we'll want to keep track of the beetle's \emph{polar angle}. This is the angle the segment $\overline{OB}$ from the center of the circular path to the beetle makes with the positive $x$-axis. The angle is measured positively in the counterclockwise direction. For now, we'll measure the polar angle in revolutions. Here's an example.

\begin{example} \label{Ex:88df83mmaz}
Between noon and 1:00pm a beetle crawls counterclockwise around a circle of radius $5$cm at a constant speed. It passes the point $(0,5)$ at 12:23pm and (without completing a full revolution) passes the point $(0,-5)$ at 12:30pm.

\begin{onlineOnly}
    \begin{center}
\desmos{zdwkilkten}{900}{600}
\end{center}
\end{onlineOnly}

\href{https://www.desmos.com/calculator/zdwkilkten}{142: Beetle on a Circle}


\begin{enumerate}
\item Find the beetle's polar angle (measured in revolutions) at 12:51pm.

\item Find a function 
\[
 \theta =a(t) \, , \, 0\leq t \leq 60, 
\]
that expresses the beetle's polar angle (in revolutions) in terms of the number of minutes past noon.
\end{enumerate}




\begin{explanation}
\begin{enumerate}
\item We start by find the beetle's signed rotation rate about the circle's center. Because the beetle's speed is constant, so is this rotation rate. The rate is positive since the beetle turns counterclockwise about the center of its path. 

At 12:23pm, when the beetle has coordinates $(0,5)$, the polar angle is
\[
 \theta = a(23) = \frac{1}{4}\text{ rev}
\] 
and at 12:30pm when the beetl has coordinates $(0,-5)$, the polar angle is
\[
  \theta = a(30) = \frac{3}{4}\text{ rev}.
\]
So the beetle rotation rate about the center of its path is
\begin{align*}
   \omega &= \frac{\Delta \theta}{\Delta t} \\
               &= \frac{(a(30) - a(23))\text{ rev}}{(30 - 23)\text{ min}} \\
               &= \frac{\left( \frac{3}{4} - \frac{1}{4}\right) \text{rev}}{7 \text{ min}} \\
               &= \frac{1}{14} \frac{\text{rev}}{\text{min}}.
\end{align*}

So the beetle's polar angle at 12:51pm is
\begin{align*}
   a(51) &= a(23) + (a(51) - a(23))\\
           &=\frac{1}{4} \text{ rev} + \left(  \frac{1}{14} \, \frac{\text{rev}}{\text{min}} \right) \left( 51-23 \right) \text{min}  \\
            &=\frac{1}{4} \text{ rev} + \left(  \frac{1}{14} \, \frac{\text{rev}}{\text{min}} \right) \left( 28 \text{ min} \right)  \\
            &= \frac{1}{4} \text{ rev} + 2 \text{ rev} \\
            &= 2.25 \text{ rev}.
\end{align*}

\item The beetle's polar angle (in revolutions) at time $t$ minutes past noon is
\begin{align*}
  a(t) &= a(23) + (a(t) - a(23))\\
           &=\frac{1}{4} \text{ rev} + \left(  \frac{1}{14} \, \frac{\text{rev}}{\text{min}} \right) \left( t-23 \right) \text{min}  \\
            &=\frac{1}{4} \text{ rev} +  \frac{1}{14}  \left( t-23\right) \text{ rev} .
           \end{align*}

Omitting units and including the domain, the polar angle function is
\[
  \theta = a(t) = \frac{1}{4} + \frac{1}{14}  \left( t-23\right) \, , \, 0\leq t\leq 60.
\]
\end{enumerate}

\end{explanation}

\end{example}




\end{document}